% =============================================================================
% Meta-Informationen der Arbeit
% -----------------------------------------------------------------------------
% HIER trägst du als Erstes deine persönlichen Daten ein. Diese Werte tauchen
% automatisch auf dem Deckblatt und in den PDF-Eigenschaften auf. Du musst sonst
% nirgends im Projekt etwas an deinen Daten ändern.
%
% Tipp: Ersetze nur den Text INNERHALB der geschweiften Klammern { ... }.
% =============================================================================

% --- Titel der Arbeit ---
\newcommand{\myTitel}{Mustertitel der wissenschaftlichen Arbeit}
\newcommand{\myUntertitel}{Ein aussagekräftiger Untertitel (optional)}

% --- Autor ---
\newcommand{\myAutor}{Vorname Nachname}
\newcommand{\myMatrikelnummer}{000000}

% --- Studiengang ---
\newcommand{\myStudiengang}{Wirtschaftsinformatik}
\newcommand{\myStudienort}{Musterstadt}

% --- Betreuung ---
\newcommand{\myBetreuer}{Prof.~Dr.~Mustername}
\newcommand{\myZweitbetreuer}{} % leer lassen, wenn nicht vorhanden

% --- Datum ---
\newcommand{\myAbgabedatum}{01.01.2026}

% --- Art der Arbeit ---
% z.B. Hausarbeit, Seminararbeit, Projektarbeit, Bachelorarbeit, Masterarbeit
\newcommand{\myArbeit}{Seminararbeit}

% --- Wortzahl (optional) ---
% Manche Pruefende verlangen die Wortzahl auf dem Titelblatt. Leer lassen,
% wenn nicht gefordert -- dann erscheint die Zeile nicht. Zaehlen so, wie
% deine Pruefenden zaehlen: ./skripte/woerter.sh <erste> <letzte Seite>
\newcommand{\myWortzahl}{}

% --- Modul / Kurs ---
\newcommand{\myModul}{Modulname}

% --- Hochschule ---
\newcommand{\myHochschule}{FOM Hochschule für Oekonomie \& Management}

% --- Firma (nur bei Arbeiten mit Sperrvermerk) ---
% Leer lassen, wenn deine Arbeit keinen Sperrvermerk braucht.
\newcommand{\myFirma}{Musterfirma GmbH}
